% !PRINTED 612
\weberpage{612}

Wenn also in einer Gruppe $L$ eine Substitution $\lambda$ vorkommt, in der $a = 1$ und $b$ von Null verschieden ist, so enthält hiernach $L$ die ganze \gesperrt{cyklische Gruppe} $(z, z + b)$ $(b = 0, 1, \ldots n - 1)$.

Ist $a$ nicht gleich 1, so setzen wir
\[
1 + a + \cdots + a^{h-1} = \frac{a^h - 1}{a - 1},
\]
und schliessen daraus, dass, wenn $e$ der kleinste positive Exponent ist, für den $a^e \equiv 1\ (\mathrm{mod}\ n)$ ist, $\lambda$ vom Grade $e$ ist.

Die Periode der Substitution $\lambda$ ist, wenn nicht $a = 1$ ist, eine intransitive Gruppe. Es giebt eine Ziffer $z$, die durch $\lambda$ und seine Wiederholungen nicht geändert wird, die aus der Congruenz $z \equiv az + b\ (\mathrm{mod}\ n)$ bestimmt wird, und mit
\[
z_0 \equiv \frac{b}{a - b}\ (\mathrm{mod}\ n)
\]
% [sic] source denominator is "a-b"; fixed-point algebra would give "1-a"
bezeichnet werden kann, wenn unter $\dfrac{\mu}{\nu}\ (\mathrm{mod}\ n)$ eine ganze Zahl verstanden wird, die durch Multiplication mit dem Factor $\nu$ nach dem Modul $n$ mit $\mu$ congruent wird.

Die cyklische Gruppe $(z, z + b)$ ist transitiv und vom Grade $n$. Es lässt sich nachweisen, dass jede transitive lineare Gruppe $(\mathrm{mod}\ n)$ die cyklische Gruppe enthalten muss.

Es sei $g$ eine primitive Wurzel von $n$ und $\alpha$ der Index von $a$, d.\,h. $g^\alpha \equiv a\ (\mathrm{mod}\ n)$ (\S.\,136). Der Kürze wegen nennen wir für den Augenblick $\alpha$ zugleich den \gesperrt{Index der Substitution} $\lambda = (z, az + b)$. Der Index kann kleiner als $n - 1$ und gleich oder grösser als Null angenommen werden. Ist er gleich Null, so ist $\lambda$ entweder die identische Substitution, oder sie gehört der cyklischen Gruppe an. Nach (4) gilt die Regel, dass der Index einer zusammengesetzten Substitution gleich der Summe der Indices der Componenten ist.

Daraus folgt, dass alle Indices der Substitutionen einer linearen Gruppe $L$ Vielfache des kleinsten positiven unter ihnen sind, und dass folglich, wenn $\alpha_0$ dieser kleinste positive Index ist und $a^{\alpha_0} = a_0$ gesetzt wird, alle Substitutionen von $L$ in der Form
% [sic] source prints a^{\alpha_0}=a_0 here; the preceding definition uses g^{\alpha}=a. Retained source-faithfully.
\[
\lambda = (z, a_0^h z + b)
\]
dargestellt werden können, worin $a_0$ festgehalten wird, und nur $h$ und $b$ gewisse Zahlenreihen durchlaufen.

Eine Substitution von der Form $\lambda_0 = (z, a_0 z + b_0)$ muss in der Gruppe $L$ vorkommen.
